\subsection{Переменные и типы данных}

Переменная — это именованная область памяти, в которой хранится значение. В Python не нужно заранее объявлять тип переменной, он определяется автоматически при присваивании.

\subsubsection*{Основные типы данных}
\begin{itemize}
    \item \texttt{int} (integer) — целые числа (например, \texttt{10}, \texttt{-5}).
    \item \texttt{float} — числа с плавающей точкой (например, \texttt{3.14}, \texttt{1.5}).
    \item \texttt{str} (string) — строки текста (например, \texttt{"Pioneer"}).
    \item \texttt{bool} (boolean) — логические значения (\texttt{True}, \texttt{False}).
\end{itemize}

\begin{codefile}[python]{python/examples/02_variables.py}
\end{codefile}

Функция \texttt{type()} позволяет узнать тип переменной.

\begin{taskbox}[Личная карточка]
Создайте переменные для хранения вашего имени, возраста и любимого предмета. Выведите их на экран одной строкой.
\end{taskbox}

\subsubsection*{Задачи для самостоятельного решения}

\begin{enumerate}
    \item Создайте переменную \texttt{speed} и присвойте ей значение 5. Выведите её на экран.
    \item Создайте переменную \texttt{height} и присвойте ей значение 1.5. Выведите её.
    \item Создайте переменную \texttt{name} со значением \texttt{"Pioneer"}. Выведите её.
    \item Создайте переменную \texttt{score} = 10. Затем измените её значение на 20. Выведите \texttt{score}.
    \item Создайте переменную \texttt{x} = 100. Затем присвойте ей строку \texttt{"Сто"}. Выведите тип переменной до и после изменения.
    \item Даны \texttt{a = 5} и \texttt{b = 7}. Создайте переменную \texttt{c}, равную их сумме. Выведите \texttt{c}.
    \item Даны \texttt{part1 = "Geo"} и \texttt{part2 = "skan"}. Создайте \texttt{full}, сложив их.
    \item Создайте строку \texttt{s = "Fly"}. Выведите её 3 раза подряд (используйте умножение).
    \item Какого типа переменная \texttt{pi = 3.1415}? Выведите результат функции \texttt{type(pi)}.
    \item Присвойте переменным \texttt{x}, \texttt{y}, \texttt{z} значения 1, 2, 3 в одной строке.
    \item Даны \texttt{a = 10}, \texttt{b = 20}. Поменяйте их значения местами без использования третьей переменной.
    \item Создайте переменную \texttt{is\_flying} со значением \texttt{True}. Выведите её.
    \item Дано целое число \texttt{n = 5}. Преобразуйте его в дробное и сохраните в \texttt{f}.
    \item Дана строка \texttt{s\_num = "123"}. Преобразуйте её в число \texttt{num} и прибавьте 1.
    \item Даны \texttt{mass = 10} и \texttt{accel = 9.8}. Вычислите \texttt{force = mass * accel}.
\end{enumerate}
